% A free particle in an expanding universe, written as a PHY 317 challenge problem (2026-11-03).
% Statement in Taylor's style with Will-style hints; solution in the course key style (inline per part,
% results boxed, the substitution shown, "Assess" at the end). Companion to AstroLagrange.tex, Sections 3-4.
%   latexmk -pdf RedshiftProblem.tex                                   -> statement + solution
%   latexmk -pdf -jobname="RedshiftProblem (statement)" -usepretex='\def\statementonly{1}' RedshiftProblem.tex
\documentclass[11pt,letterpaper]{article}
\usepackage[margin=0.75in]{geometry}
\usepackage{amsmath,amssymb,mathtools}
\usepackage{xcolor}
\usepackage[most]{tcolorbox}
\usepackage{enumitem}
\usepackage{parskip}
\usepackage[colorlinks=true, linkcolor=blue, citecolor=blue, urlcolor=blue]{hyperref}

\newif\ifsolutions \solutionstrue
\ifdefined\statementonly\solutionsfalse\fi
\NewDocumentCommand{\taylor}{m}{\begin{tcolorbox}[enhanced,colback=gray!4,colframe=gray!45,boxrule=0.5pt,
    left=3pt,right=3pt,top=3pt,bottom=3pt]#1\end{tcolorbox}}
\setlist[enumerate]{label=(\alph*), leftmargin=2em, itemsep=0.4em}
\newcommand{\Lag}{\mathcal{L}}
\newcommand{\Ham}{\mathcal{H}}
\newcommand{\vpec}{v_{\rm pec}}
\newcommand{\pd}[2]{\frac{\partial #1}{\partial #2}}

\begin{document}
\noindent\textbf{PHYSICS 317: a challenge problem on Lagrange's equations}\ifsolutions\ --- \textbf{solutions}\fi

\taylor{\textbf{Challenge problem (optional): a free particle in an expanding universe.}
In Newtonian cosmology, space is filled with matter of uniform density $\rho(t)$, and every
distance grows by one common factor: a particle that just rides the expansion sits at
$r = a(t)\,x$ with $x$ fixed, where $a(t)$ is the \emph{scale factor} and $x$ is its
\emph{comoving coordinate}. By Newton's shell theorem the matter inside radius $r$ pulls a
particle there as if all of it sat at the origin (and the matter outside exerts no net force),
so a particle of mass $m$ at distance $r$ from the origin has potential energy
\[
  U(r, t) = \frac{2\pi G}{3}\,m\rho(t)\,r^2
  \qquad\Bigl[\text{check: } -\pd{U}{r} = -\frac{Gm}{r^2}\cdot\frac{4\pi}{3}\rho r^3 = -\frac{GmM(<r)}{r^2}\Bigr],
\]
and the scale factor obeys the Friedmann equation $\ddot a = -\frac{4\pi G}{3}\rho\,a$ (Newton's
second law for a particle riding the expansion, with $r = ax$; you may take it as given).
Now consider a particle of mass $m$ that moves along a line through the origin but does
\emph{not} necessarily ride the expansion.
\begin{enumerate}
\item Write down the Lagrangian $\Lag(r, \dot r, t)$. Find the Lagrange equation for $r$ and
  check that it is Newton's second law with the force above. Explain why neither of our two
  conservation tests (an ignorable coordinate, Section~7.6; no explicit $t$ in $\Lag$,
  Section~7.8) hands you a conserved quantity here.
\item Change to the comoving coordinate $x = r/a(t)$ (a generalized coordinate whose relation
  to $r$ involves $t$, as in Example~7.6). Show that
  \[
    \Lag = \tfrac12 m a^2\dot x^2 + \frac{d}{dt}\Bigl(\tfrac12 m a\dot a\,x^2\Bigr).
  \]
  \emph{Hint:} expand the square in $\tfrac12 m\dot r^2$; write the cross term
  $m a\dot a\,x\dot x$ as a total time derivative minus what is left over; use the Friedmann
  equation on what is left over.
\item Explain why the total time derivative can be dropped without changing the equation of
  motion (Problem~7.48, or argue directly from Hamilton's principle). From here on take
  $\Lag = \tfrac12 m a(t)^2\dot x^2$.
\item Which coordinate is now ignorable? Find the conserved generalized momentum, and show that
  the \emph{peculiar velocity} $\vpec \equiv a\dot x$, the particle's velocity relative to the
  expansion, falls as $1/a$. (When the universe doubles in size, a galaxy's random motion halves.)
\item Find the Hamiltonian $\Ham = p\dot x - \Lag$. Is it $T + U$? Is it conserved? Use
  $d\Ham/dt = -\partial\Lag/\partial t$ (Section~7.8) to find how $\Ham$ depends on $a$.
\item General relativity gives the same conservation law for light: the momentum a local
  observer measures is the conserved momentum divided by $a$. With $p_{\rm local} = h/\lambda$,
  deduce the cosmological redshift $1 + z \equiv \lambda_{\rm obs}/\lambda_{\rm emit} =
  a_{\rm obs}/a_{\rm emit}$.
\end{enumerate}}

\ifsolutions
\textbf{(a)} $\Lag = T - U$ with $T = \tfrac12 m\dot r^2$:
\[
  \boxed{\;\Lag(r, \dot r, t) = \tfrac12 m\dot r^2 - \frac{2\pi G}{3}\,m\rho(t)\,r^2\;}
\]
Euler-Lagrange time! $\partial\Lag/\partial\dot r = m\dot r$ and
$\partial\Lag/\partial r = -\frac{4\pi G}{3} m\rho\, r$, so
$\frac{d}{dt}\partial\Lag/\partial\dot r = \partial\Lag/\partial r$ reads
\[
  \boxed{\; m\ddot r = -\frac{4\pi G}{3}\,m\rho\, r \;}
\]
which is $m\ddot r = -GmM(<r)/r^2$ with $M(<r) = \frac{4\pi}{3}\rho r^3$: Newton's second law
with the shell-theorem force. Neither test applies: $r$ appears in $\Lag$ (through $U$), so it
is not ignorable, and $t$ appears explicitly (through $\rho(t)$), so Section~7.8's guarantee of
a constant $\Ham$ is off. (A conserved quantity does exist, as (d) shows; these coordinates just
do not reveal it.)

\textbf{(b)} With $r = ax$, the product rule gives $\dot r = \dot a x + a\dot x$. Substituting
into (a),
\[
  \Lag = \tfrac12 m(\dot a x + a\dot x)^2 - \frac{2\pi G}{3}\, m\rho\, a^2 x^2
  = \tfrac12 m a^2\dot x^2 + m a\dot a\, x\dot x + \tfrac12 m\dot a^2 x^2 - \frac{2\pi G}{3}\, m\rho\, a^2 x^2 .
\]
The cross term: by the product rule (and the chain rule for $\frac{d}{dt}x^2 = 2x\dot x$),
\[
  \frac{d}{dt}\Bigl(\tfrac12 m a\dot a\,x^2\Bigr)
  = \tfrac12 m(\dot a^2 + a\ddot a)\,x^2 + m a\dot a\,x\dot x
  \quad\Longrightarrow\quad
  m a\dot a\,x\dot x = \frac{d}{dt}\Bigl(\tfrac12 m a\dot a\,x^2\Bigr) - \tfrac12 m\dot a^2 x^2 - \tfrac12 m a\ddot a\,x^2 .
\]
Put that in:
\[
  \Lag = \tfrac12 m a^2\dot x^2 + \frac{d}{dt}\Bigl(\tfrac12 m a\dot a\,x^2\Bigr)
  - \tfrac12 m\dot a^2 x^2 - \tfrac12 m a\ddot a\,x^2 + \tfrac12 m\dot a^2 x^2 - \frac{2\pi G}{3}\, m\rho\, a^2 x^2 .
\]
The two $\tfrac12 m\dot a^2 x^2$ terms cancel. Friedmann, $\ddot a = -\frac{4\pi G}{3}\rho\,a$,
turns the $a\ddot a$ term into
$-\tfrac12 m a\bigl(-\frac{4\pi G}{3}\rho\,a\bigr)x^2 = +\frac{2\pi G}{3}\,m\rho\,a^2x^2$, which
cancels the potential term exactly. What is left is
\[
  \boxed{\;\Lag = \tfrac12 m a^2\dot x^2 + \frac{d}{dt}\Bigl(\tfrac12 m a\dot a\,x^2\Bigr)\;}
\]
(The whole of $U$ went into the $a\ddot a$ term: the expansion is \emph{caused} by that
gravity, so in coordinates that expand with it, the gravity is already accounted for.)

\textbf{(c)} Write $F = \tfrac12 m a\dot a\,x^2$, a function of $x$ and $t$ only. Adding $dF/dt$
to $\Lag$ adds $\int_{t_1}^{t_2}(dF/dt)\,dt = F(t_2) - F(t_1)$ to the action, and Hamilton's
principle holds $x(t_1)$ and $x(t_2)$ fixed, so that number is the same for every path being
compared. The stationary path, and so the equation of motion, is unchanged. (Problem~7.48 checks
the same thing at the level of Lagrange's equations: $\partial(dF/dt)/\partial x$ and
$\frac{d}{dt}\partial(dF/dt)/\partial\dot x$ are both $\partial_x\partial_t F + \dot x\,\partial_x^2 F$,
so they cancel.)

\textbf{(d)} With $\Lag = \tfrac12 m a^2\dot x^2$, $x$ does not appear: $x$ is ignorable, and
\[
  \boxed{\; p \equiv \pd{\Lag}{\dot x} = m a^2\dot x \text{ is conserved} \;}
\]
since Lagrange's equation reads $\frac{d}{dt}(m a^2\dot x) = \partial\Lag/\partial x = 0$.
In terms of the peculiar velocity, $p = m a\,(a\dot x) = m a\,\vpec$, and $p$ is constant, so
\[
  \boxed{\;\vpec \propto \frac1a\;}
\]
Double the universe, halve every peculiar velocity.

\textbf{(e)} With $p = m a^2\dot x$,
\[
  \Ham = p\dot x - \Lag = m a^2\dot x^2 - \tfrac12 m a^2\dot x^2 = \tfrac12 m a^2\dot x^2
  = \boxed{\;\tfrac12 m\vpec^2\;}
\]
the kinetic energy in the peculiar velocity. It is not $T + U$: $T = \tfrac12 m\dot r^2 =
\tfrac12 m(\dot a x + a\dot x)^2$ contains the Hubble-flow part, and $U$ is not in $\Ham$ at all.
That is allowed, because the relation $r = a(t)x$ between the generalized coordinate and the
position involves $t$ (Section~7.8's condition for $\Ham = T + U$ fails). It is not conserved
either: $\Lag$ depends on $t$ through $a$, so by the chain rule
$\partial\Lag/\partial t = \tfrac12 m(2a\dot a)\dot x^2 = m a\dot a\,\dot x^2$, and
\[
  \frac{d\Ham}{dt} = -\pd{\Lag}{t} = -m a\dot a\,\dot x^2 = -2\,\frac{\dot a}{a}\cdot\tfrac12 m a^2\dot x^2
  = -2\,\frac{\dot a}{a}\,\Ham
  \quad\Longrightarrow\quad \boxed{\;\Ham \propto \frac1{a^2}\;}
\]
($d\ln\Ham/dt = -2\,d\ln a/dt$ integrates to $\ln\Ham = -2\ln a + \text{const}$.) That agrees
with (d): $\Ham = \tfrac12 m\vpec^2$ and $\vpec \propto 1/a$.

\textbf{(f)} The local momentum is $p/a$ with $p$ constant, so $p_{\rm local} \propto 1/a$. For
a photon $p_{\rm local} = h/\lambda$, so $h/\lambda \propto 1/a$, i.e.\ $\lambda \propto a$, and
\[
  \boxed{\;1 + z = \frac{\lambda_{\rm obs}}{\lambda_{\rm emit}} = \frac{a_{\rm obs}}{a_{\rm emit}}\;}
\]
A photon emitted when the universe was half its present size arrives with twice its wavelength:
$z = 1$.

\textbf{Assess:} The two results are the same statement about the same symmetry. In comoving
coordinates nothing depends on $x$, so the momentum conjugate to $x$ is conserved; measured
locally that momentum falls as $1/a$, which is Hubble drag for a galaxy and the redshift for a
photon. Energy is not conserved, and the formalism says exactly why: $t$ is in $\Lag$, the
expanding universe has no time-translation symmetry, and $d\Ham/dt = -\partial\Lag/\partial t$
gives the rate. The one hand-wave is the shell theorem on an infinite medium; general relativity
makes that honest and gives (f) without borrowing.
\fi
\end{document}
